\section{Discussion}\label{sec:disc}

\noindent\textbf{What Coordination Changes in Practice.} The empirical evidence demonstrates that introducing stateful, bounded-authority agents genuinely alters the behavioral dynamics of the interview preparation system. By tracking conceptual weaknesses across turns, the Strategy Agent smooths difficulty transitions, significantly reducing high-frequency level oscillations ($-2.67$) and directional reversals ($-1.88$) compared to the frozen single-policy PPO baseline. However, scientific rigor demands that we distinguish behavioral stabilization from pedagogical quality. Replay on recorded candidate answers cannot determine whether a smoother difficulty progression promotes superior candidate learning, retention, or interview confidence. A smoother trajectory demonstrates algorithmic consistency, not verified educational superiority.

\noindent\textbf{The Limits and Realities of Agent Specialization.} A pervasive narrative in the multi-agent literature assumes that decomposing complex reasoning into specialized agents inherently produces superior task performance~\cite{autogen,metagpt}. Our findings directly challenge this assumption in deterministic settings. When agents operate over deterministic heuristic rules and equivalent information, specialization produces \textbf{zero algorithmic advantage}: decision outputs and difficulty trajectories are identical to an ablated single-agent or unified configuration ($M_3 = M_1$, 160/160 identical sessions). 

Instead, the defensible justification for agent specialization is strictly architectural:
\begin{enumerate}
\item \emph{Decoupled Lifecycle Engineering:} Strategy and Coaching address fundamentally distinct concerns (macro-level progression vs micro-level pedagogical presentation). Specialization enables independent rule authoring, targeted unit testing, and isolated feature flagging without risking unintended regressions across domains.
\item \emph{Bulkhead Fault Isolation:} Distributing state across private ledgers prevents failure propagation. An unhandled exception or malformed payload in Coaching leaves Strategy difficulty adaptation fully operational, whereas a monolithic controller drops both functions.
\item \emph{Operational Overhead:} This modularity is not free; it incurs a measurable runtime cost ($+1.04$\,ms median latency per turn; 47.7 invocations per session) due to repeated serialization, Pydantic schema validation, and message dispatch.
\end{enumerate}

\noindent\textbf{Bounded Authority vs Emergent Multi-Agent Intelligence.} Our architecture intentionally eschews open-ended agent negotiation or decentralized consensus mechanisms. In high-stakes educational and evaluation settings, unconstrained agent communication introduces risks of specification drift, circular reasoning, and cascading hallucination~\cite{mast}. By enforcing a centralized orchestrator as the sole system of record, restricting agents to typed proposal contracts, and applying strict threshold-gated arbitration, our system realizes the principle of least privilege~\cite{ss75}. The system achieves predictable, verifiable coordination without sacrificing centralized safety bounds.

\noindent\textbf{Relevance to ADCoM 2027.} This investigation aligns directly with the core themes of the 31st International Conference on Advanced Computing and Communications (ADCoM 2027), particularly Track 2: \emph{Agent-Native Development and Lifecycle Engineering}. By rigorously auditing contract boundaries, quantifying the trade-offs of agent specialization, and isolating coordination mechanisms through cluster-bootstrap replay, this work demonstrates how multi-agent architectures can be engineered, bounded, and evaluated with empirical integrity.
